\fichatitulo{54 · METODOLOGIA}{Journal of Management Information Systems · 2007}{Design Science Research Methodology}
{\small PEFFERS, Ken; TUUNANEN, Tuure; ROTHENBERGER, Marcus A.; CHATTERJEE, Samir. \textit{A Design Science Research Methodology for Information Systems Research}. Journal of Management Information Systems, v. 24, n. 3, p. 45--77, 2007--2008. \href{https://doi.org/10.2753/MIS0742-1222240302}{DOI} e PDF consultado.\par}

\campo{Problema e objetivo}
Como oferecer um processo reconhecível para conduzir, apresentar e avaliar pesquisas de Design Science em sistemas de informação.

\campo{Principal contribuição · síntese interpretativa}
Propõe a DSRM como modelo nominal de seis atividades: identificação e motivação do problema; definição dos objetivos da solução; projeto e desenvolvimento; demonstração; avaliação; e comunicação. A metodologia também funciona como modelo mental para relatar e avaliar artefatos.

\campo{Método / natureza da evidência}
Artigo metodológico que fundamenta o processo em literatura anterior e o demonstra por quatro estudos de caso de artefatos de sistemas de informação.

\campo{Resultado ou argumento principal}
A aplicação aos quatro casos mostra como um artefato pode ser ligado ao problema, aos objetivos, à demonstração e à avaliação. A DSRM organiza o processo, mas não elimina a necessidade de justificar o método de avaliação conforme o artefato e o contexto.

\campo{Excertos-chave · seleção literal}
\begin{itemize}
\excerto{Processo}{problem identification and motivation}{Resumo, p. 45.}
\excerto{Atividades}{design and development, demonstration, evaluation, and communication}{Resumo, p. 45.}
\end{itemize}

\campo{Limitações e alcance}
\textbf{Fonte:} o artigo apresenta um modelo nominal e quatro demonstrações, não um protocolo estatístico universal.\textbf{Análise da ficha:} a DSRM sustenta a organização do artefato avaliativo, mas não prova a eficácia do protocolo da dissertação.

\campo{Aproveitamento na dissertação · proposta}
\textbf{Destino:} metodologia. \textbf{Uso:} relacionar problema parlamentar, artefato de avaliação, demonstração, avaliação empírica e comunicação dos resultados.

\registro{Fonte consultada: PDF local, 34 páginas; resumo, introdução, atividades da DSRM, demonstração, avaliação e conclusão. A edição impressa aparece como Winter 2007--8; a chave bibliográfica original foi preservada. Chave: \texttt{peffersEtAl2007dsrm}.}
